\hypertarget{persistence}{}
\hypertarget{rolling-persistence-overlap-and-quintile-mobility}{%
\chapter{Rolling persistence, overlap and quintile
mobility}\label{rolling-persistence-overlap-and-quintile-mobility}}

\hypertarget{the-overlap-illusion}{%
\section{The overlap illusion}\label{the-overlap-illusion}}

Two 120-month windows ending 12 months apart share 108 months, or 90\%
of their data. High correlation between their rankings is therefore
partly mechanical. A 120-month separation produces disjoint decades and
is the structurally interpretable comparison.

\[
Overlap fraction = max(0, W - h) / W
\]

For W = 120 and horizon h = 12, the fraction is 0.90; for h = 120, it is
zero.

\hypertarget{persistence-metrics}{%
\section{Persistence metrics}\label{persistence-metrics}}

\begin{itemize}
\tightlist
\item
  \textbf{Spearman rank correlation:} agreement in order, invariant to
  monotone score transformations.
\item
  \textbf{Pearson score correlation:} linear agreement in estimated
  posterior scores.
\item
  \textbf{Mean absolute rank change:} average number of positions moved.
\item
  \textbf{Transition matrix:} probability of moving from one quintile to
  another.
\item
  \textbf{Jaccard overlap:} intersection divided by union for top/bottom
  membership sets.
\end{itemize}

\[
J(A,B) = |A \cap  B| / |A \cup  B|
\]

\hypertarget{correct-interpretation}{%
\section{Correct interpretation}\label{correct-interpretation}}

\begin{longtable}[]{@{}p{0.14\textwidth}p{0.22\textwidth}p{0.22\textwidth}p{0.30\textwidth}@{}}
\toprule\noalign{}
Horizon & Window overlap & Mean rank correlation & Interpretation \\
\midrule\noalign{}
\endhead
\bottomrule\noalign{}
\endlastfoot
12 months & 90\% & 0.863 & Mostly an overlap diagnostic \\
60 months & 50\% & 0.483 & Still contaminated by overlap \\
120 months & 0\% & 0.156 & Weak structural persistence \\
\end{longtable}

The same-quintile probability at the disjoint 120-month horizon is about
25\%, only modestly above the 20\% uniform-mobility benchmark. This
warns against treating the point ordering as a stable hierarchy.

